\subsection{包络代数的定义}

设 g 是 K 上的李代数。对于 K 上任意含单位元的结合代数 L，从 g 到 L 的 \(\alpha\)-映射是从 g 到 L 的 K-线性映射 \(\sigma\)，满足

\[
\sigma ([ x, y ]) = \sigma (x) \sigma (y) - \sigma (y) \sigma (x) \quad (x, y \text {   in   } g)
\]

（换言之，是从 \(\mathfrak{g}\) 到与 L 相关联的李代数的同态）。若 L' 是 K 上另一个含单位元的结合代数，\(\tau\) 是将 1 映到 1 的从 L 到 L' 的同态，则 \(\tau \circ \sigma\) 是一个 \(\alpha\)-映射，从 \(\mathfrak{g}\) 到 L'。我们要寻找一个含单位元的结合代数以及一个 \(\alpha\)-映射，从 \(\mathfrak{g}\) 到该代数，二者具有泛性质（Set Theory, Chapter IV, § 3, no. 1）。

定义 1. 设 \(\mathfrak{g}\) 是 \(\mathbf{K}\) 上的李代数，\(\mathbf{T}\) 是 \(\mathbf{K}\)-模 \(\mathfrak{g}\) 的张量代数，\(\mathbf{J}\) 是 \(\mathbf{T}\) 的双边理想，由张量 \(x \otimes y - y \otimes x - [x, y]\)（其中 \(x \in \mathfrak{g}, y \in \mathfrak{g}\)）生成。结合代数 \(\mathbf{U} = \mathbf{T} / \mathbf{J}\) 称为 \(\mathfrak{g}\) 的包络代数。在 \(\mathfrak{g}\) 上，将 \(\mathbf{T}\) 映到 \(\mathbf{U}\) 的典范映射的限制称为从 \(\mathfrak{g}\) 到 \(\mathbf{U}\) 的典范映射。

令 \(T_{+}\) 为 T 的双边理想，由 0 阶分量为零的张量组成。令 \(\mathrm{T}_0 = \mathrm{K}.1\) 为 T 中 0 阶元素的集合。令 \(\mathrm{U}_{+}\) 和 \(\mathrm{U}_0\) 为 \(\mathrm{T}_{+}\) 和 \(\mathrm{T}_0\) 在 U 中的典范像。由于 \(\mathrm{J} \subset \mathrm{T}_{+}\)，直和分解 \(\mathrm{T} = \mathrm{T}_0 + \mathrm{T}_{+}\) 蕴含直和分解 \(\mathrm{U} = \mathrm{U}_0 + \mathrm{U}_{+}\)。因此代数 U 有一个异于 0 的单位元，且 \(\mathrm{U}_0 = \mathrm{K}.1\)。对所有 \(x \in \mathbf{U}\)，\(x\) 在 \(\mathrm{U}_0\) 中的分量称为 \(x\) 的常数项。常数项为零的元素构成 U 的双边理想，即由 g 在 U 中的典范像生成的双边理想 \(\mathrm{U}^{+}\)。

结合代数 U 由 1 和 g 在 U 中的典范像生成。

若 \(x \in \mathfrak{g}\) 且 \(y \in \mathfrak{g}\)，则 \(x \otimes y - y \otimes x\) 与 \([x, y]\) 在 T 中模 J 同余；因此，若 \(\sigma_0\) 表示从 \(\mathfrak{g}\) 到 U 的典范映射，则

\[
\sigma_ {0} (x) \sigma_ {0} (y) - \sigma_ {0} (y) \sigma_ {0} (x) = \sigma_ {0} ([ x, y ])
\]

在 U 中成立。换言之，\(\sigma_{0}\) 是从 g 到 U 的 \(\alpha\)-映射。

命题 1. 设 \(\sigma\) 是一个 \(\alpha\)-映射，从 \(\mathfrak{g}\) 到含单位元的结合代数 \(\mathbf{L}\)。存在唯一的同态 \(\tau\)，它从 \(\mathbf{U}\) 到 \(\mathbf{L}\)，将 1 映到 1，且满足 \(\sigma = \tau \circ \sigma_0\)，其中 \(\sigma_0\) 表示从 \(\mathfrak{g}\) 到 \(\mathbf{U}\) 的典范映射。

设 \(\tau'\) 是从 T 到 L 的唯一同态，它延拓 \(\sigma\) 并将 1 映到 1。则对 g 中的 \(x, y\)，

\[
\tau^ {\prime} (x \otimes y - y \otimes x - [ x, y ]) = \sigma (x) \sigma (y) - \sigma (y) \sigma (x) - \sigma ([ x, y ]) = 0;
\]

故 \(\tau'\) 在 \(J\) 上为零，取商后定义了同态 \(\tau\)，它从 \(U\) 到 \(L\)，将 1 映到 1，并满足 \(\sigma = \tau \circ \sigma_0\)。由于 \(\tau\) 由 \(\sigma_0(g)\) 和 1 在代数 \(U\) 中生成，其唯一性是显然的。

设 \(g'\) 是 K 上的另一个李代数，U' 是其包络代数，\(\sigma_0'\) 是从 \(g'\) 到 U' 的典范映射。设 \(\phi\) 是从 g 到 g' 的同态。则 \(\sigma_0' \circ \phi\) 是从 g 到 U' 的 \(\alpha\)-映射；因此存在唯一的将 1 映到 1 的从 U 到 U' 的同态 \(\tilde{\phi}\)，使下图

\[
\begin{array}{c} g \xrightarrow {\phi} g ^ {\prime} \\ \sigma_ {0} \Biggl \downarrow \\ U \xrightarrow {\tilde {\phi}} U ^ {\prime} \end{array}
\]

是交换的。此同态将 U 中常数项为零的元素映到 U' 中常数项为零的元素。若 g'' 是 K 上的另一个李代数，\(\phi'\) 是从 g' 到 g'' 的同态，则 \(\widetilde{\phi'\circ\phi}=\widetilde{\phi}'\circ\widetilde{\phi}\)。

